\section{研究方案}

\subsection{研究内容}
\begin{enumerate}[label=\arabic*., leftmargin=2em]
    \item \textbf{三维框架（杆、梁）结构有限元离散方案：}针对冷弯薄壁型钢框架结构中杆、梁构件的几何与受力特点，研究适用的低阶/高阶单元列式及减缩积分、假设应变等免自锁策略，建立静力、动力及稳定分析统一框架，并发展等效外荷载计算与节点校核方案。
    \item \textbf{梁杆结构静动力分析测试与验证：}设计覆盖轴向、弯曲、剪切及组合受力状态的算例，结合理论解、商业有限元软件（ABAQUS/ANSYS）结果以及课题组自主程序 ApproxOperator.jl 的计算结果，对模块的精度、收敛性与计算效率进行系统验证。
    \item \textbf{多国规范校核方案制定：}针对中国 GB 50018《冷弯薄壁型钢结构技术规范》、北美 AISI S100、欧洲 EC3 等主流规范，建立截面强度、构件稳定、整体稳定及连接验算的接口与流程，实现分析结果向规范校核的自动化映射。
    \item \textbf{模块集成与接口开发：}将有限元分析模块与 BuildingMaster（BM）设计软件及生产控制系统进行数据对接，支持几何模型、荷载、材料与计算结果的自动化导入导出，形成“设计—分析—校核—生产”的闭环。
    \item \textbf{工程案例应用与反馈优化：}依托大禾众邦提供的典型轻钢建筑项目，开展模块的工程应用测试，根据用户反馈持续优化计算流程、结果可视化与交互体验。
\end{enumerate}

\subsection{技术路线}
本项目的技术路线如图~\ref{fig:tech-route} 所示。首先，通过文献调研与企业需求访谈，明确冷弯薄壁型钢框架结构有限元分析的功能需求与性能指标；其次，基于课题组自主程序 ApproxOperator.jl，发展免自锁的杆/梁单元列式，构建静力、动力与稳定分析的统一计算框架；再次，设计覆盖典型受力状态的测试算例，与理论解及 ABAQUS/ANSYS 结果进行对比验证；随后，针对 GB 50018、AISI S100、EC3 等规范建立校核接口；最后，将模块集成至 BuildingMaster（BM）设计软件，并依托大禾众邦工程案例开展应用测试与迭代优化。

\begin{figure}[htbp]
    \centering
    \fbox{\parbox{0.8\textwidth}{\centering \vspace{2cm} 技术路线图（请替换为实际图片）\vspace{2cm}}}
    \caption{项目技术路线}
    \label{fig:tech-route}
\end{figure}

\subsection{研究方法}
本项目采用文献调研、理论推导、数值实验、对比验证与工程应用相结合的研究方法：

\begin{itemize}[leftmargin=2em]
    \item \textbf{文献调研法：}系统梳理冷弯薄壁型钢结构有限元分析、免自锁单元、稳定分析及相关规范校核方面的国内外研究进展，明确技术路线与关键科学问题。
    \item \textbf{理论推导法：}基于变分原理与有限元方法，推导适用于薄壁杆/梁结构的单元刚度矩阵、质量矩阵及几何/材料非线性列式，建立免自锁分析方案。
    \item \textbf{数值实验法：}在 ApproxOperator.jl 平台上实现算法原型，通过标准算例测试单元的收敛性、精度与计算效率。
    \item \textbf{对比验证法：}将本项目计算结果与理论解、ABAQUS/ANSYS 商业软件结果及试验数据进行对比，验证模块的可靠性。
    \item \textbf{工程应用法：}依托大禾众邦提供的真实工程案例，开展模块应用测试，收集用户反馈并持续优化。
\end{itemize}

评价指标主要包括：位移/应力计算精度、屈曲临界荷载误差、计算收敛阶、单工况计算耗时、与商业软件结果的一致性以及规范校核覆盖率。

